\subsection{Berechnung der Ausgangsspannung bei zwei Eingangsspannungen}
% Alt: Aufgabe 4
\Aufgabe{
    \label{Aufg4}
    Gegeben ist die abgebildete Schaltung mit einem idealen
    Operationsverstärker. Dabei ist der Widerstand $R_\mathrm{1} = \mathrm{100\,k\Omega}$ und $R_2=\mathrm{10\,k\Omega}$. Die Betriebsspannung beträgt $\mathrm{\pm 15\,V}$.
    \begin{itemize}
        \item [ a)] Ermitteln Sie die Ausgangsspannung $U_\mathrm{A}$ bei $U_\mathrm{E}=2\,\mathrm{mV}$ und $U_\mathrm{bat} = \mathrm{0\,V}$
        \item [ b)] Ermitteln Sie die Ausgangsspannung $U_\mathrm{A}$ bei $U_\mathrm{E}=20\,\mathrm{mV}$ und $U_\mathrm{bat}=\mathrm{0\,V}$
    \end{itemize}
    \begin{figure}[H]
        \centering
        \input{Tikz/src/FigOPV2Eingaenge.tex}\alttext{Das Bild zeigt ein elektrisches Schaltbild mit zwei Widerständen, beschriftet als \(R_1\) und \(R_2\), die in Reihe geschaltet sind. Zwischen den Widerständen ist ein Operationsverstärker mit unendlicher Verstärkung dargestellt. Der Operationsverstärker hat einen Plus- und einen Minus-Eingang, wobei der Plus-Eingang mit einer Batterie verbunden ist, die mit \(U_{\mathrm{bat}}\) beschriftet ist. Die Spannungen an verschiedenen Punkten im Schaltkreis sind mit Pfeilen und den Bezeichnungen \(U_\mathrm{E}\), \(U_\mathrm{X}\) und \(U_\mathrm{A}\) markiert. Alle Spannungsquellen sind jeweils mit Masse verbunden, dargestellt durch das Erdungssymbol.}
        \label{fig:FigOPV2Eingaenge}
    \end{figure}

}

\Loesung{
    %\begin {figure} [H]
    %\centering
    %\input{Tikz/src/LsgOPV2Eingaenge.tex}
    %\label{fig:LsgOPV2Eingaenge}
    %\end {figure}
    Verstärkung berechnen:
    \begin{align*}
        V =\mathrm{- \frac{100\,k\Omega}{10\,k\Omega} = -10}
    \end{align*}
    Formel für Verstärkung:
    \begin{align*}
        V = \frac{U_\mathrm{A}}{U_\mathrm{E}}
    \end{align*}
    Formel umstellen nach $U_\mathrm{A}$:
    \begin{align*}
        U_\mathrm{A} = V \cdot U_\mathrm{E}
    \end{align*}
    \begin{itemize}
        \item [ \bf a)]
              \begin{align*}
                  U_\mathrm{A} = \mathrm{-10 \cdot 2\,mV = -20\,mV}
              \end{align*}
        \item [ \bf b)]
              \begin{align*}
                  U_\mathrm{A} = \mathrm{-10 \cdot 20\,mV = -200\,mV}
              \end{align*}
    \end{itemize}
    Siehe: Abschnitt \ref{fig: Verschaltung Verstaerker} und Tabelle \ref{tab:Grundschaltungen}
}